\documentclass{sintefbeamer}
\usepackage[utf8]{inputenc}
\usepackage[T1]{fontenc}
\usepackage{geometry}
\usepackage{graphicx}
\usepackage{float}
\usepackage{hyperref}
\usepackage[english,indonesian]{babel}

\title{Optimasi Pemrosesan Template LaTeX Menggunakan Pendekatan Berbasis AI}
\author{Academic Support Agent}
\date{\today}

\begin{document}
\maketitle

\begin{abstract}

\end{abstract}

\chapter{PENDAHULUAN}


\chapter{TINJAUAN PUSTAKA}


% Original content (commented out):
% \documentclass{sintefbeamer}
\newcommand{\testcolor}[1]{\colorbox{#1}{\textcolor{#1}{test}}~\texttt{#1}}
% meta-data
\titlebackground{images/fondo}

\title{FEN Beamer\\ Tema para Presentaciones}
\subtitle{Usando \LaTeX\ para crear diapositivas}
\author{\href{mailto:cquirozv@fen.uchile.cl}{César Quiroz V.}}
\date{\today}

% No borrar latexmkrc, dedido a que es el archivo que deja la fecha y hora de Chile
% document body
\begin{document}

\maketitle

\begin{frame}

Esta plantilla es una adaptación de \hrefcol{https://www.overleaf.com/latex/templates/harbin-institute-of-technology-hit-beamer-presentation-theme/prwxqwfdzkqj}{(HIT) Beamer Presentation Theme} creada por \hrefcol{mailto:qilong-kirov.liu@connect.hit.cn}{LIU Qi-long (Kirov)}, con contenidos de \hrefcol{https://www.overleaf.com/latex/templates/statale-unimi-presentation/ykkwvfdbqydr}{Statale Unimi Presentation} creada por \href{mailto:ciao@gio.im}{Giorgio Marchetti}. La plantilla original en la que se han basado ellos es \hrefcol{https://www.overleaf.com/latex/templates/sintef-presentation/jhbhdffczpnx}{SINTEF Presentation} creada por \hrefcol{mailto:federico.zenith@sintef.no}{Federico Zenith} \vspace{\baselineskip}

A continuación hay una pequeña introducción escrita por \hrefcol{mailto:federico.zenith@sintef.no}{Federico Zenith} sobre como usar \LaTeX\ y beamer para crear diapositivas. Todos los derechos reservados por él\vspace{\baselineskip}

This template is released under \hrefcol{https://creativecommons.org/licenses/by-nc/4.0/legalcode}{Creative Commons CC BY 4.0} license

\end{frame}

\section{Introducción}

\begin{frame}{Beamer for SINTEF slides}{\thesection \, \secname}

\begin{itemize}
\item Asumimos que sabes utilizar \LaTeX; si no es así,
\hrefcol{http://en.wikibooks.org/wiki/LaTeX/}{puedes aprenderlo aquí}
\item Beamer es uno de los tipos de documentos (document
classes) más populares y poderosos para crear presentaciones en \LaTeX
\item Beamer también tiene un detallado
\hrefcol{http://www.ctan.org/tex-archive/macros/latex/contrib/beamer/doc/beameruserguide.pdf}{manual de usuario}
\item Aquí presentaremos solo las funciones más básicas para que te puedas poner al día
\end{itemize}
\end{frame}

\begin{frame}{Beamer vs. PowerPoint}
Comparado con PowerPoint, usar \LaTeX\ es mejor debido a que:
\begin{itemize}
\item Esto no es lo que ves es lo que obtienes, sino que es \emph{lo que quieres decir} es lo que obtienes:\\
escribes el contenido, el computador hace la composición (typesetting).
\item Crea un \texttt{pdf}: no hay problemas con fuentes, fórmulas,
      versiones de programas, etc.
\item Es más fácil mantener estilos, fuentes, resaltados, etc. consistentes.
\item La escritura de fórmulas en \TeX\ es la mejor:
\begin{equation*}
\mathrm{i}\,\hslash\frac{\partial}{\partial t} \Psi(\mathbf{r},t) =
-\frac{\hslash^2}{2\,m}\nabla^2\Psi(\mathbf{r},t)
+ V(\mathbf{r})\Psi(\mathbf{r},t)
\end{equation*}

\end{itemize}
\end{frame}

\section{Edición}

\begin{frame}[fragile]{Selección tipo de documento}
Después de la última actualización del perfil gráfico, el tema \texttt{sintef} para
Beamer ha sido actualizado completamente a un tipo de documento.
Para comenzar a trabajar con \texttt{sintefbeamer}, crea un documento de \LaTeX\ con el
preámbulo:
\begin{block}{Código mínimo documento SINTEF Beamer }
\begin{lstlisting}[language=TeX]
\documentclass{sintefbeamer}
\begin{document}
\begin{frame}{Hello, world!}
\end {frame}
\end {document}
\end{lstlisting}
\end{block}
\end{frame}

\begin{frame}[fragile]{Portada}
Para crear una portada típica, tienes que escribir los siguientes comandos:
\begin{block}{Comandos para la Portada}
\begin{lstlisting}[language=TeX]
\title{Sample Title}
\subtitle{Sample subtitle}
\author{First Author, Second Author}
\date{Defaults to today's}
\end{lstlisting}
\end{block}
Después puedes escribir \verb|\maketitle| para crear la diapositiva.

Para cambiar la imagen de fondo tienes que usar el comando
\verb|\titlebackground| antes de \verb|\maketitle|; su único argumento es el nombre del archivo (y/o la ubicación).

\end{frame}

\begin{frame}[fragile]{Escribiendo una diapositiva simple}
\framesubtitle{It's really easy!}
\begin{itemize}
\item Una diapositiva típica tiene listas con viñetas
\item[1)] La viñeta se puede cambiar modificando lo que está dentro de los corchetes [ ]
\end{itemize}
\begin{block}{Código para Diapositiva con listado}<+->
\begin{lstlisting}[language=TeX]
\begin{frame}
  \frametitle{Writing a Simple Slide}
  \framesubtitle{It's really easy!}
  \begin{itemize}
    \item A typical slide has bulleted lists
    \item These can be uncovered in sequence
  \end{itemize}
\end{frame}\end{lstlisting}
\end{block}
\end{frame}

\begin{frame}[fragile]{Agregar Imágenes}
\begin{columns}
\begin{column}{0.7\textwidth}
Agregar imágenes se hace de la misma forma que en otro documento \LaTeX:
\begin{block}{Código para agregar imágenes}
\begin{lstlisting}[language=TeX]
\usepackage{graphicx}
% ...
\includegraphics
[width=\textwidth]{images/default}
\end{lstlisting}
\end{block}
\end{column}
\begin{column}{0.3\textwidth}
\includegraphics
[width=\textwidth]{images/default}\\
\end{column}
\end{columns}
\end{frame}

\begin{frame}[fragile]{Dividir en Columnas}
Dividir la diapositiva es fácil;
comúnmente, un lado contiene una imagen y el otro el texto:
\begin{columns}
\begin{column}{0.6\textwidth}
Esta es la primera columna
\end{column}
\begin{column}{0.3\textwidth}
Y esta es la segunda
\end{column}
\end{columns}
\begin{block}{Column Code}
\begin{lstlisting}[language=TeX]
\begin{columns}
    \begin{column}{0.6\textwidth}
        This is the first column
    \end{column}
    \begin{column}{0.3\textwidth}
        And this the second
    \end{column}
    % There could be more!
\end{columns}
\end{lstlisting}
\end{block}
\end{frame}

\begin{frame}[fragile]
\frametitle{Fuente}
\begin{itemize}
\item La tarea principal de las fuentes es que se puedan leer.
\item Aquí hay algunas buenas...
  \begin{itemize}
  \item {\textrm{Usa fuentes serif solo con proyectores de alta definición}}
  \item {\textsf{De lo contrario (o si lo prefieres) usa fuentes sans-serif }}
  \end{itemize}
\item ... y otras no muy buenas:
  \begin{itemize}
  \item {\texttt{Nunca uses monospace para texto normal}}
  \item {\frakfamily Gothic, calligraphic o fuentes raras: intenta evitarlas siempre}
\end{itemize}
\end{itemize}
\end{frame}

\section{Personalización}

\footlinecolor{RUChile}
\begin{frame}[fragile]{Cambiando los Estilos de las Diapositivas}
\begin{itemize}
\item Puedes elegir blanco o el \textit{color principal} de la diapositiva \emph{en el preámbulo} with \verb|\themecolor{white}| (por defecto) or \verb|\themecolor{blue}|
      \begin{itemize}
      \item \emph{No} deberías cambiar esto a lo largo del documento: A Beamer no le gusta
      \item Si \emph{realmente lo debes hacer} tienes que agregar 
      \verb|\usebeamercolor[fg]{normal text}| en la diapositiva
      \end{itemize}
\item Puedes cambiar el \textbf{color del pie de página} con 
\verb|\footlinecolor{color}|
      \begin{itemize}
      \item Escribe el comando \emph{antes} de una nueva \verb|frame|
      \item Hay cuatro colores ``oficiales'': 
      \testcolor{sintefblue} (color principal), \testcolor{AoUChile}, 
      \testcolor{RUChile}, \testcolor{VUChile}
      \item Por defecto no hay pie de página; lo puedes agregar con 
      \verb|\footlinecolor{}|
      \item Otros colores pueden funcionar, pero no está garantizado.
      \end{itemize}
\end{itemize}
\end{frame}

\begin{frame}[fragile]{Cuadros}
\begin{columns} % agregando [onlytextwidth] los márgenes izquierdos serán establecidos correctamente.
\begin{column}{0.3\textwidth}
\begin{block}{Cuadros estándar}
Estos están coordinados con el color del pie de página (y son grises en el caso del tema azul)
\begin{verbatim}
\begin{block}{title}
content...
\end{block}
\end{verbatim}
\end{block}
\end{column}
\begin{column}{0.7\textwidth}
\begin{colorblock}[black]{sinteflightgreen}{Cuadros de color}
Son similares al de la izquierda, pero tú eliges el color. El texto es blanco por defecto, pero lo puedes cambiar con un argumento adicional.
\small
\begin{verbatim}
\begin{colorblock}[black]{sinteflightgreen}{title}
content...
\end{colorblock}
\end{verbatim}
\end{colorblock}
Algunos ejemplos de colores son: \testcolor{sinteflilla}, 
\testcolor{sintefblue}, \testcolor{sintefdarkgreen}, 
\testcolor{sintefyellow}, \testcolor{airforceblue}
\end{column}
\end{columns}
\end{frame}

\footlinecolor{}
\begin{frame}[fragile]{Usando los Colores}
\begin{itemize}[<alert@2>]
  \item Puedes usar los colores con el comando
        \verb|\textcolor{<color name>}{text}| 
  \item Los colores están definidos en el paquete \texttt{sintefcolor}:
  \begin{itemize}
  \item Colores primarios: \testcolor{sintefblue} y su compañero 
  \testcolor{sintefgrey}
  \item Tres tonos de verde: \testcolor{sinteflightgreen}, 
  \testcolor{sintefgreen}, \testcolor{sintefdarkgreen}
  \item Y otros adicionales: \testcolor{sintefyellow}, \testcolor{sintefred}, 
        \testcolor{sinteflilla}
        \begin{itemize}
        \item Puedes cambiar los tonos---revisa la información de \verb|sintefcolor| o el  \hrefcol{https://sintef.sharepoint.com/sites/stottetjenester/%
        kommunikasjon/grafisk-profil-new/Sider/default.aspx}{manual de SINTEF}
        \end{itemize}
  \end{itemize}
  \item \emph{No} abuses de los colores: \verb|\emph{}| es usualmente suficiente.
  \item Usa \verb|\alert{}| para generar \alert<2->{atención} en algún lado.
  \item<2- | alert@2> Si tu resaltas mucho, ¡al final no resaltaste nada!
\end{itemize}
\end{frame}

\begin{frame}[fragile]{Otros}
\begin{itemize}
\item Para cambiar entre los temas claro y oscuro, utiliza
      \texttt{light} (por defecto) o \texttt{dark}. No es posible cambiar el tema solo para una diapositiva.
\item Para agregar una diapositiva final usa \verb|\backmatter|.
\begin{itemize}
      \item El título también aparece en el pie de página junto al autor, puedes cambiar esto con \verb|\footlinepayoff|
      \item Puedes remover el título de la diapositiva final con \verb|\backmatter[notitle]|
      \end{itemize}
\item El aspecto está definido en 16:9, pero lo puedes cambiar a 4:3 para proyectores antiguos agregando \texttt{aspectratio=43} en 
\verb|\documentclass[]{}|; cualquier otro valor aceptado por Beamer tambien es posible.
\end{itemize}
\end{frame}

\begin{chapter}[images/Fondo]{sintefblue}{Diapositivas especiales}
\begin{itemize}
\item Diapositivas por secciones
\item Diapositivas con imágenes laterales
\end{itemize}
\end{chapter}

\footlinecolor{sintefblue}
\begin{frame}[fragile]{Diapositivas por Secciones}
\begin{itemize}
\item Son similares a \verb|frame|s, pero con algunas opciones adicionales
\item Se crean con \verb|\begin{chapter}[<image>]{<color>}{<title>}|
\item La imagen es opcional, el color y el título son obligatorios.
\item Hay 7 colores ``oficiales'': \testcolor{sintefblue}, 
\testcolor{sintefdarkgreen}, \testcolor{sintefgreen}, 
\testcolor{sinteflightgreen}, \testcolor{sintefred}, \testcolor{sintefyellow}, 
\testcolor{sinteflilla}.
      \begin{itemize}
      \item Extrañamente, hay \emph{más} colores para el pie de página.
      \item Puede ser un buen toque cambiar el color del pie de página acorde al color del tema. Tú decides.
      \end{itemize}
\item De todas formas, \verb|chapter| se comporta de la misma forma que \verb|frame|.
\end{itemize}
\end{frame}

\begin{sidepic}{images/default}{Diapositivas con Imágenes Laterales}
\begin{itemize}
\item Se crean con \texttt{$\backslash$begin\{sidepic\}\{<image>\}\{<title>\}}
\item De lo contrario, \texttt{sidepic} trabaja de igual forma que \texttt{frame}
\end{itemize}
\end{sidepic}

\backmatter

\end{document}


\end{document}